\subsection{二元集合公理}

\[
(\forall x) (\forall y) \operatorname{Coll} _ {z} (z = x \text {   或   } z = y).
\]

此公理说明，若 \(x\) 并且 \(y\) 是对象，则存在一个集合，其唯一元素为 \(x\) 并且 \(y\) .

定义 2. 集合 \(\mathfrak{E}_z\) ( \(z = x\) 或 \(z = y\) ），其唯一元素为 \(x\) 和 \(y\) ，记作 \(\{x, y\}\) .

关系 \(z \in \{x, y\}\) 因此等价于“ \(z = x\) 或 \(z = y\) ”；由 C50 推出 \(\{y, x\} = \{x, y\}\) .

让 \(R\{z\}\) 为一个关系，并设 \(x, y\) 为不同于 \(z\) 的字母。由准则 C32、C33（第一章，§4，第 3 号）和 C43（第一章，§5，第 1 号）容易推出关系 \((\exists z)((z \in \{x, y\})\) 并且 \(R\{z\})\) 等价于“ \(R\{x\}\) 要么 \(R\{y\}\) ”；因此关系 \((\forall z)((z \in \{x, y\}) \Longrightarrow R\{z\})\) 等价于“ \(R\{x\}\) 并且 \(R\{y\}\) ''.

集合 \(\{x, x\}\) ，简记为 \(\{x\}\) ，称为唯一元素为 x 的集合；关系 \(z \in \{x\}\) 等价于 z = x，且关系 \(x \in X\) 等价于 \(\{x\} \subset X\) .

